% Les cinq images cachées : double page finale (verso + recto en vis-à-vis).
% Les images viennent du manuel 2026 (code/easter_egg_images.py de ce manuel), copiées dans figures/.
% Ne jamais écrire ici ce que les noms représentent pour la famille.
\newcommand{\oeuf}[4]{% #1 nom, #2 fichier, #3 phrase, #4 développement de l'acronyme
  \noindent{\sffamily\bfseries\LARGE\color{accent}\textls[180]{#1}}\hfill
    {\footnotesize\itshape\color{encre!70}#4}\par\vspace{4pt}
  \noindent\begin{tcolorbox}[enhanced,colback=white,colframe=encre!12,boxrule=0.5pt,arc=2pt,
      left=3pt,right=3pt,top=3pt,bottom=3pt,halign=center,drop fuzzy shadow=encre!25]
    \includegraphics[width=\linewidth,height=3.95cm,keepaspectratio]{#2}\end{tcolorbox}\par\vspace{2pt}
  \noindent{\small #3}\par}

\cleartoverso
\xglobal\colorlet{accent}{okorange}
\phantomsection
\addcontentsline{toc}{chapter}{Les cinq images cachées}
\markboth{Les cinq images cachées}{Les cinq images cachées}
\thispagestyle{empty}
\vspace*{-0.6cm}
{\fontsize{25}{30}\selectfont Les cinq images cachées\par}
\vspace{6pt}{\color{okorange}\rule{2.6cm}{1.4pt}}\par\vspace{8pt}
\noindent{\small\itshape Les quatre lettres d'ALIX et les trois de KAL se cachent dans ces théorèmes…}\par
\vspace{4pt}
\noindent{\footnotesize\color{encre!80}Une chasse au trésor pour finir. Cinq théorèmes portent des noms courts ; chacun
redonne, sous ce nom, un théorème déjà vérifié ailleurs par le juge, et ils sont rangés tous les cinq dans le fichier
\texttt{Acronymes.lean}. Les théorèmes d'origine, rejoués le 26 octobre, sont de niveau~A (pour ALIX, sous réserve du théorème amont
d'openai/math) ; les cinq noms courts en sont des alias exacts, non rejoués eux-mêmes. Les images sont des calculs de niveau~X, qui illustrent un cas particulier sans rien démontrer.}\par
\vspace{14pt}
\oeuf{ALIX}{egg_ALIX.pdf}{À n'importe quelle densité, le réseau triangulaire a la plus petite énergie possible pour toute
force répulsive de la bonne famille, si le théorème amont d'openai/math est juste.}{A Lattice Is eXtremal}
\vspace{16pt}
\oeuf{KAL}{egg_KAL.pdf}{La chance qu'une mutation finisse par envahir toute la population est une limite : celle de la
probabilité d'avoir été « absorbé » par l'état où tous la portent, quand le nombre de générations grandit.}%
{Kolmogorov Absorption Limit}
\clearpage

\thispagestyle{empty}
\vspace*{-0.6cm}
\oeuf{ILENE}{egg_ILENE.pdf}{Sur un réseau où chacun est remplacé au même rythme, un mutant a la même chance de
l'emporter, quel que soit l'endroit où il apparaît : c'est la loi isotherme.}{Isothermal Law: Every Node Equivalent}
\vspace{9pt}
\oeuf{ALEXIS}{egg_ALEXIS.pdf}{Un mélange positif d'exponentielles décroissantes garde, pour toutes ses dérivées, les
signes alternés qui font une force admissible : les mélanges de Laplace gardent leurs signes.}%
{Averaged Laplace EXponentials Inherit Signs}
\vspace{9pt}
\oeuf{AURELIA}{egg_AURELIA.pdf}{Sous la sélection, quand chaque type garde sa fitness, la fitness moyenne de la
population ne diminue jamais.}{Average Under Replication: Evolution Lifts It, Always}
\clearpage
